\documentclass[12pt,a4paper]{article}
\usepackage[utf8]{inputenc}
\usepackage[T1]{fontenc}
\usepackage{amsmath,amssymb,amsthm}
\usepackage{geometry}
\usepackage{hyperref}
\usepackage{booktabs}

\geometry{margin=1in}

\newtheorem{theorem}{Theorem}
\newtheorem{axiom}{Axiom}
\newtheorem{proposition}{Proposition}
\newtheorem{lemma}{Lemma}
\newtheorem{conjecture}{Conjecture}

\theoremstyle{definition}
\newtheorem{definition}{Definition}

\theoremstyle{remark}
\newtheorem{remark}{Remark}

\title{\textbf{Co-emergence of mass, signature, time, and Hilbert space}\\[6pt]
{\large DRAFT}}

\author{Will Regelmann\thanks{Corresponding author.} \quad Claude (Anthropic)\thanks{AI
assistant; contributed to mathematical development and analysis.}}

\date{March 2026}

\begin{document}
\maketitle

\begin{abstract}
Lorentzian self-consistency produces complex phases in the fixed-point
wavefunction that are absent in the Riemannian case. These phases force
subsystem entanglement entropy to exceed the Riemannian baseline (proved
for rank-2 subsystems; conjectured for general rank, with zero violations
in over $20{,}000$ random matrices through dimension $16 \times 16$). A finite toy model
confirms: composition, Born rule statistics, and imaginary subsystem
coherences all emerge from Lorentzian self-consistency through $N = 16$,
with phase structure
stable through $N = 256$. The Lorentzian and Riemannian fixed points
share identical magnitude profiles by construction (the curvature map
depends only on $|\psi_\sigma|^2$), so the $68\%$ entropy excess is
entirely a phase effect.

We propose that this mechanism extends to a co-emergence of mass,
Lorentzian signature, local time, and local Hilbert space: four aspects
of a single self-consistent structure, none existing without the others.
Physical mass requires hyperbolicity (Wigner 1939); hyperbolicity
provides local proper time along timelike geodesics; local proper time
enables a local Hilbert space via the Page--Wootters mechanism; and the
resulting quantum stress-energy sources gravity, closing the
self-consistency loop. We formalize this co-emergence within a four-level
hierarchy of geometric constraints on a block 4-manifold with no evolution
parameter, no background structure, and no fundamental Hilbert space.
The semiclassical Einstein equation --- the hierarchy's metric level ---
is signature-blind: self-consistent solutions exist for both Lorentzian
and Riemannian metrics, for both massive and massless field content.
Signature selection operates not at any single level but at the
cross-level constraint, where only Lorentzian metrics with massive field
content support the full hierarchy. The mass assumption itself can be
weakened: the conformal trace anomaly generates effective mass for
non-conformally coupled fields, closing the co-emergence loop without
bare mass as input.

\bigskip\noindent\textit{Note:} Companion to ``Fixed-point existence of
self-consistent semiclassical gravity''~\cite{fixed_point}, whose results
are cited but not re-derived.
\end{abstract}

%======================================================================
\section{Introduction}
\label{sec:intro}
%======================================================================

Every formulation of fundamental physics assumes Lorentzian signature.
The distinction between one timelike and three spacelike directions---the
$(-,+,+,+)$ sign pattern in the metric---is an input, not a derivation.
If the universe is a four-dimensional geometric block satisfying
constraints rather than a state evolving in time, this distinction should
be a consequence of the constraints, not a separate postulate.

The most natural route to deriving signature would be topological: the
intersection form $Q_\mathcal{M}$ of a 4-manifold has a type $(b^+, b^-)$
that mirrors the $(1,3)$ split of Lorentzian signature. But this route is
obstructed. The intersection form is a global invariant on
$H_2(\mathcal{M};\mathbb{Z})$, while metric signature is a pointwise
property of the tangent bundle; no theorem connects them. Worse, closed
simply-connected 4-manifolds with nontrivial intersection forms have
Euler characteristic $\chi \geq 3$, while a closed manifold admits a
Lorentzian metric only if $\chi = 0$~\cite{geroch}. The manifold class
where the Freedman--Donaldson machinery~\cite{freedman,donaldson} is
strongest cannot carry a Lorentzian metric at all.

The semiclassical Einstein equation (SCE) provides no help either. The
companion paper~\cite{fixed_point} establishes that the perturbative
self-consistency map has contraction constant $\kappa \sim (m/M_P)^2$
regardless of metric signature. The Banach contraction is
signature-blind: self-consistent solutions exist for both Lorentzian and
Riemannian metrics. The Starobinsky de Sitter solution~\cite{starobinsky}
extends this blindness to massless conformal matter. Level by level,
signature is not selected.

The foundational tension in quantum gravity is well known: general
relativity requires no preferred time parameter, while standard quantum
mechanics requires one. The Wheeler--DeWitt equation
$\hat{H}\Psi = 0$~\cite{dewitt,wheeler} is the canonical attempt at
resolution---the Hamiltonian constraint makes the wave functional
timeless. But every known formulation of $\hat{H}\Psi = 0$ introduces
hidden temporal or foliation structure at some stage of its
derivation~\cite{kuchar}. The canonical ADM approach requires a $3+1$
decomposition; the covariant path integral hides foliation dependence in
the Wick rotation, boundary conditions, and measure ambiguity; the WKB
semiclassical limit inserts time via Born--Oppenheimer
adiabaticity~\cite{halliwell,banks}.

This paper proposes that signature and time are two faces of the same
coin, and that neither can be derived in isolation. Mass, Lorentzian
signature, local time, and local Hilbert space are four aspects of a
single self-consistent structure: each requires the others, and all
emerge together from the requirement that the geometric block satisfy
constraints at every level simultaneously. We call this the
\emph{co-emergence thesis}. It is formalized as a four-level hierarchy
of geometric constraints (Section~\ref{sec:hierarchy}) whose joint
solution space is restricted to Lorentzian signature with massive field
content (Conjecture~\ref{conj:coemergence}). Moreover, the mass
assumption itself can be weakened: the conformal trace anomaly generates
effective mass for non-conformal fields on any self-consistent
background with nonzero curvature
(Conjecture~\ref{conj:mass_generation}).

The paper is organized as follows. Section~\ref{sec:axioms} states the
three axioms. Section~\ref{sec:phase} presents the mathematical results:
the phase structure of the Lorentzian fixed point, the entropy excess
lemma, and a finite toy model confirming composition, Born rule
statistics, and imaginary subsystem coherences.
Section~\ref{sec:coemergence} develops the
co-emergence thesis that these results motivate: mass, signature, time,
and Hilbert space as four links of a closed logical chain, together with
the self-consistent mass generation mechanism that weakens the mass
assumption. Section~\ref{sec:hierarchy} develops the four-level hierarchy.
Section~\ref{sec:dimensions} argues that four-dimensionality is derived,
not assumed. Section~\ref{sec:connections} connects the hierarchy to
Wheeler--DeWitt, the SCE, Page--Wootters, and Osterwalder--Schrader
reconstruction. Section~\ref{sec:open} lists priority open problems.

%======================================================================
\section{Axioms}
\label{sec:axioms}
%======================================================================

\begin{axiom}[No evolution parameter]
\label{ax:timeless}
The fundamental description contains no evolution parameter. All equations
are constraints on the structure of the four-dimensional block, not rules for
generating it. A formulation satisfies this axiom if and only if it can be
stated entirely as a constraint on a four-dimensional geometric structure with
no preferred slicing.
\end{axiom}

\begin{axiom}[No background structure]
\label{ax:nobackground}
The fundamental description contains no fixed metric, no preferred
foliation, no assumed symmetry group, and no presupposed dimensionality
or signature. The consistency constraints are purely geometric: at the
deepest level, metric structure and matter both emerge from
topology~(Conjecture).
\end{axiom}

\begin{axiom}[No fundamental Hilbert space]
\label{ax:nohilbert}
The fundamental description does not take a Hilbert space as input.
Hilbert space, if it appears, is derived---an effective description for
subsystems with limited access to the full block.
\end{axiom}

Axiom~\ref{ax:nohilbert} requires justification, since every existing
formulation of quantum gravity starts from a Hilbert space.

A Hilbert space $\mathcal{H}$ with inner product
$\langle\cdot|\cdot\rangle$ carries temporal structure in three ways,
any one of which conflicts with Axiom~\ref{ax:timeless}:

\begin{enumerate}

\item \textbf{The inner product requires a time-slice.} In quantum field
theory, $\langle\Psi_1|\Psi_2\rangle = \int \Psi_1^*[\phi]\Psi_2[\phi]
\,\mathcal{D}\phi$ integrates over field configurations $\phi$ on a spatial
hypersurface at a fixed time. Different choices of hypersurface give
inequivalent inner products and inequivalent quantum theories.

\item \textbf{Unitary evolution requires a time parameter.} The statement
$|\Psi(t_2)\rangle = U(t_2, t_1)|\Psi(t_1)\rangle$ requires a preferred
parameter $t$. Even in the ``frozen'' Wheeler--DeWitt formulation, the
Hilbert space inner product is defined on a spatial slice, smuggling in
the foliation through the back door.

\item \textbf{The Born rule requires a preferred decomposition.}
The probability $|\langle\phi|\Psi\rangle|^2$ is defined relative to a
measurement basis $\{|\phi\rangle\}$. Selecting this basis is an
observer-dependent choice that the fundamental description should not
require.

\end{enumerate}

\begin{proposition}[Hilbert space violates Axiom~\ref{ax:timeless}]
\emph{(Rigorous for (1) and (2); Sketch for (3).)}
\label{prop:nohilbert}
Any formulation of quantum gravity that takes a Hilbert space with inner
product as a fundamental input violates Axiom~\ref{ax:timeless}, because the
inner product requires a preferred spatial hypersurface (point~(1)) or a
preferred time parameter (point~(2)).
\end{proposition}

\noindent\emph{Argument.} Point~(1) follows from the standard QFT
definition of the Fock space inner product: the functional integral
$\int\mathcal{D}\phi$ is over field configurations on a Cauchy surface,
and different Cauchy surfaces give different inner products. Point~(2) is
immediate from the definition of unitary evolution. \hfill$\square$

\medskip

The conclusion is that a formulation satisfying Axiom~\ref{ax:timeless}
must not take Hilbert space as input. If Hilbert space appears at all,
it must emerge as a structure for subsystems---the content of
Section~\ref{sec:coemergence}.

\medskip

\noindent\textbf{What this paper is and is not.} This is a framework
proposal, not a completed theory. The co-emergence thesis
(Section~\ref{sec:coemergence}) rests on specific mathematical claims at
stated rigor levels. The metric level has the most rigorous support---it
is the setting of the companion paper~\cite{fixed_point}. The
topological, smooth-structure, and effective-QM levels are research
programs with clear mathematical agendas. We do not claim to have derived
the Standard Model. We claim to have identified the structure that
connects mass, signature, time, and Hilbert space, and to have made
explicit what is required to prove the connection rigorously.

%======================================================================
\section{Lorentzian phases and the entropy excess}
\label{sec:phase}
%======================================================================

The co-emergence thesis (Section~\ref{sec:coemergence}) requires that
marginalizing the smooth-structure measure $\mu^*$ over the
complement of a subregion produces quantum statistics---density matrices,
interference, Born rule probabilities. Whether this is possible depends
on a structural property of $\mu^*$: is $\mu^*$ real-valued or
complex-valued?

A real-valued positive measure cannot produce destructive interference.
If $\mu^*$ is a probability measure in the ordinary sense---real,
positive, normalized---then the marginal over the complement of
$\mathcal{R}$ is a classical conditional probability, and the effective
theory for $\mathcal{R}$ is a classical stochastic process. No amount
of marginalization or coarse-graining can generate complex amplitudes
from real-valued weights. If $\mu^*$ is real, the framework produces
statistical mechanics, not quantum mechanics.

Quantum statistics require complex weights that can cancel. In the path
integral, this structure comes from the Lorentzian action: the weight
$e^{iS/\hbar}$ is a complex phase whose oscillations produce the
interference that distinguishes quantum from classical. On a Riemannian
manifold, the corresponding weight $e^{-S_E/\hbar}$ is real and positive.

This observation connects the phase structure of $\mu^*$ directly to
signature selection. If the smooth-structure measure inherits its weight
from the action functional on each smooth structure, then:
\begin{itemize}
\item On a Lorentzian manifold, $\mu^*$ carries complex phases $\to$
marginalization can produce quantum statistics $\to$ Level~3 is
nontrivial.
\item On a Riemannian manifold, $\mu^*$ is real-valued $\to$
marginalization produces only classical correlations $\to$ Level~3
collapses to classical statistics.
\end{itemize}
This is a sharper mechanism for Level~3 signature selectivity than the
absence of timelike geodesics alone. The Riemannian fixed point fails
not only because there are no local clocks, but because the measure
itself lacks the phase structure required for quantum interference.

%----------------------------------------------------------------------
\subsection{Phase structure of the fixed point}
\label{sec:phase_structure}
%----------------------------------------------------------------------

The complex phases of $\mu^*$ are a property of the \emph{global}
fixed point, not of any single subsystem's density matrix. In the toy
model (Section~\ref{sec:toy_model}), the Lorentzian $\psi^*$ carries
$\sim 25\%$ of its norm in the imaginary part for rank-2 bipartitions
(after global phase alignment), rising to $\sim 34\%$ for higher-rank
subsystems and stable across all system sizes tested. However, when a
large environment is traced out, the off-diagonal imaginary content of
the \emph{reduced} density matrix is suppressed by partial-trace
averaging: contributions from different environment configurations
carry different phases, which partially cancel in the sum. The quantum
signature of the phases survives in the eigenvalue spectrum of the
reduced density matrix---measurably, as a $68\%$ excess in entanglement
entropy ($S_{\mathrm{Lor}} / S_{\mathrm{Riem}} \approx 1.68$, stable
across all scaling regimes). Stripping all phases from the Lorentzian
$\psi^*$ exactly reproduces the Riemannian fixed point itself, not
merely its entropy: since the curvature map $R_\sigma(\psi)$ depends
on $\psi$ only through $|\psi_\sigma|^2$, both the Lorentzian and
Riemannian fixed-point equations reduce to the same equation for the
magnitudes $|\psi^*_\sigma|$. The Lorentzian fixed point is
$|\psi^*_\sigma|^{1+i\theta}$ (up to normalization) and the Riemannian
is $|\psi^*_\sigma|$, with the same $|\psi^*_\sigma|$ in both cases.
Phase-stripping recovers the Riemannian state by construction, so the
entropy excess is entirely a phase effect.
Lemma~\ref{lem:entropy_excess} proves that this excess is guaranteed
for any non-uniform magnitude profile.

The distinction matters: a na\"ive diagnostic that checks only whether
the reduced density matrix has large imaginary off-diagonal elements
will see a signal that fades with environment dimension and may conclude
the mechanism fails. The correct diagnostic examines either the global
state or the eigenvalue structure of the reduced state. The phases are
there; they restructure the subsystem's spectrum even when the
reduced-matrix representation partially averages them away.

%----------------------------------------------------------------------
\subsection{Phase-induced entropy excess}
\label{sec:entropy_excess}
%----------------------------------------------------------------------

\begin{lemma}[Phase-induced purity decrease and entropy excess]
\textbf{(Rigorous.)}
\label{lem:entropy_excess}
Let $M$ be a $d_{\mathrm{sub}} \times d_{\mathrm{env}}$ matrix with
entries $m_{ij} > 0$, normalized so that $\|M\|_F = 1$. Define
$M(\theta)$ entrywise by $M(\theta)_{ij} = m_{ij}^{1 + i\theta}$,
renormalized to $\|M(\theta)\|_F = 1$. Since $|m^{1+i\theta}| = m$,
the Frobenius norm is invariant before renormalization, so
$M(\theta)_{ij} = m_{ij} \cdot e^{i\theta \log m_{ij}}$: the
magnitudes are unchanged and the phases are locked to the magnitudes by
$\varphi_{ij} = \theta \log m_{ij}$.

Let $\rho(\theta) = M(\theta)^\dagger M(\theta)$ be the Gram matrix
(reduced density matrix), $\{\sigma_k\}$ the singular values of $M$
and $\{\sigma_k(\theta)\}$ those of $M(\theta)$, with $S = -\sum_k
\sigma_k^2 \log \sigma_k^2$ the entanglement entropy. Then for all
$\theta \neq 0$:
\begin{enumerate}
\item[(a)] \textbf{(All ranks.)} The purity decreases:
$\mathrm{Tr}(\rho(\theta)^2) \leq \mathrm{Tr}(\rho(0)^2)$,
equivalently $S_2(\theta) \geq S_2(0)$ where $S_2 = -\log
\mathrm{Tr}(\rho^2)$ is the R\'enyi-2 entropy.
\item[(b)] \textbf{(Rank 2.)} The von Neumann entropy increases:
$S(\theta) \geq S(0)$.
\end{enumerate}
In both cases, equality holds if and only if all $m_{ij}$ are equal.
\end{lemma}

\begin{proof}
We prove three structural facts and derive (a) and (b).

\textbf{Fact 1} (Frobenius norm preserved)\textbf{:} $\sum_k
\sigma_k(\theta)^2 = \sum_k \sigma_k^2 = 1$.

\textbf{Fact 2} (Largest singular value decreases)\textbf{:}
$\sigma_1(\theta) \leq \sigma_1$. By definition,
\[
  \sigma_1(\theta) = \max_{\|x\| = \|y\| = 1} |y^\dagger M(\theta)\, x|
  = \max_{\|x\| = \|y\| = 1} \Big|\sum_{ij} \bar{y}_i\, m_{ij}\,
  e^{i\theta \log m_{ij}}\, x_j \Big|.
\]
By the triangle inequality,
\[
  \Big|\sum_{ij} \bar{y}_i\, m_{ij}\, e^{i\theta \log m_{ij}}\, x_j
  \Big| \leq \sum_{ij} |y_i|\, m_{ij}\, |x_j|
  \leq \max_{\|u\| = \|v\| = 1} \sum_{ij} |v_i|\, m_{ij}\, |u_j|
  = \sigma_1,
\]
where the last equality uses the Perron--Frobenius property: for a
matrix with non-negative entries, the largest singular value is achieved
by non-negative singular vectors.

\textbf{Fact 3} (Entrywise Gram contraction)\textbf{:} The diagonal of
$\rho(\theta)$ is preserved, $\rho(\theta)_{jj} = \rho(0)_{jj}$, and
the off-diagonal entries contract: $|\rho(\theta)_{jk}| \leq
\rho(0)_{jk}$ for all $j \neq k$. By direct computation,
\[
  \rho(\theta)_{jk} = \sum_i m_{ij}\, m_{ik}\,
  e^{i\theta(\log m_{ik} - \log m_{ij})}.
\]
For $j = k$ the phases vanish, giving $\rho(\theta)_{jj} = \sum_i
m_{ij}^2 = \rho(0)_{jj}$. For $j \neq k$, the triangle inequality
gives $|\rho(\theta)_{jk}| \leq \sum_i m_{ij}\, m_{ik} =
\rho(0)_{jk}$.

\textbf{Part (a):} From Fact~3,
\[
  \mathrm{Tr}(\rho(\theta)^2) = \sum_{j,k} |\rho(\theta)_{jk}|^2
  = \sum_j \rho(0)_{jj}^2 + \sum_{j \neq k} |\rho(\theta)_{jk}|^2
  \leq \sum_j \rho(0)_{jj}^2 + \sum_{j \neq k} \rho(0)_{jk}^2
  = \mathrm{Tr}(\rho(0)^2).
\]
Equality requires $|\rho(\theta)_{jk}| = \rho(0)_{jk}$ for all
$j \neq k$, which forces each triangle inequality in Fact~3 to be
tight---all phases $\theta(\log m_{ik} - \log m_{ij})$ must be equal
across $i$ for each pair $(j,k)$, hence all $m_{ij}$ equal.

\textbf{Part (b):} When $\min(d_{\mathrm{sub}}, d_{\mathrm{env}}) =
2$, there are at most two nonzero singular values $\sigma_1 \geq
\sigma_2 \geq 0$ with $\sigma_1^2 + \sigma_2^2 = 1$. Since
$\sigma_1(\theta) \leq \sigma_1$ (Fact~2) and $\sigma_1(\theta)^2 +
\sigma_2(\theta)^2 = 1$ (Fact~1), it follows that $\sigma_2(\theta)
\geq \sigma_2$. The squared singular values are more uniform, so
$S(\theta) \geq S(0)$. The equality condition follows as in Part~(a).
\end{proof}

\begin{remark}[Von Neumann entropy for general rank]
\label{rem:vn_general_rank}
For $\min(d_{\mathrm{sub}}, d_{\mathrm{env}}) \geq 3$, the von Neumann
entropy excess $S(\theta) \geq S(0)$ does \emph{not} hold for all
positive matrices. A counterexample: the $3 \times 3$ matrix with
entries spanning a ratio of $\sim\!10^7$ has $S(\theta) < S(0)$ for
$\theta \gtrsim 5$, and matrices with sufficiently extreme entry
heterogeneity can violate $S(\theta) \geq S(0)$ at any $\theta > 0$.
The mechanism is that while $\sigma_1(\theta)$ decreases, the remaining
singular values can redistribute non-uniformly---in particular, the
smallest singular value can collapse toward zero, reducing entropy
despite the top singular value decreasing.

For matrices with moderate entry heterogeneity (entries drawn from
log-normal distributions with unit variance in log-space, corresponding
to entry ratios $\lesssim e^3 \approx 20$), the von Neumann excess has
been verified numerically for over $100{,}000$ matrices at dimensions
up to $16 \times 16$ and $\theta$ up to $20$, with zero violations.
The toy model entries $m_\sigma = e^{-R_\sigma}$ fall in this regime.
\end{remark}

\begin{remark}[Application to the toy model]
\label{rem:entropy_application}
In the toy model, the Lorentzian fixed point
$\psi^*_\sigma \propto \exp((-1 + i\theta)\, R_\sigma)$ has magnitude
$|\psi^*_\sigma| \propto e^{-R_\sigma}$ and phase $\theta R_\sigma$,
which is exactly the structure $m^{1+i\theta}$ of
Lemma~\ref{lem:entropy_excess} with $m_\sigma = e^{-R_\sigma}$. The
Riemannian fixed point has $\theta = 0$, giving the same magnitudes
but no phases.

Lemma~\ref{lem:entropy_excess}(a) rigorously guarantees that the
Lorentzian state has lower purity (higher R\'enyi-2 entropy) than the
Riemannian state for all bipartitions, whenever the curvature values
$R_\sigma$ are not all equal---i.e., whenever the external field $h$
is non-uniform. For the von Neumann entropy: the toy model entries
$m_\sigma = e^{-R_\sigma}$ with $h \sim \mathrm{Uniform}[0.5, 1.5]$
have moderate heterogeneity (entry ratio $\lesssim 3$), placing them
well within the regime where $S(\theta) \geq S(0)$ holds
(Remark~\ref{rem:vn_general_rank}). For rank-2 bipartitions,
Lemma~\ref{lem:entropy_excess}(b) proves the von Neumann excess
rigorously.

The entropy ratio depends on $\theta$ (quadratically
for small $\theta$: $S(\theta)/S(0) \approx 1 + c\,\theta^2$) and on
the magnitude heterogeneity controlled by~$h$. For $\theta = 1$ and
$h \sim \mathrm{Uniform}[0.5, 1.5]$, numerics give
$S_{\mathrm{Lor}}/S_{\mathrm{Riem}} \approx 1.69$, stable across all
system sizes tested.
\end{remark}

%----------------------------------------------------------------------
\subsection{Finite toy model}
\label{sec:toy_model}
%----------------------------------------------------------------------

A finite truncation ($N = 4$, $\mathrm{Sm}(\mathcal{M}) = S_R \times
S_C$ with $|S_R| = |S_C| = 2$) with an explicit self-consistency map
$\mathcal{F}$ partially answers whether self-consistency naturally
produces complex-valued fixed points and whether the marginal has the
algebraic structure of a density matrix. In the minimal model:
\begin{itemize}
\item The Lorentzian coupling $\gamma = -1 + i\theta$ produces
complex-valued fixed points whose configuration-dependent relative
phases survive the partial trace, yielding a subsystem density matrix
$\rho_R$ with complex off-diagonal elements ($|\mathrm{Im}\,
\rho_{01}| \sim 10^{-2}$ at $\theta = 1$). The phases originate
in the oscillatory weight $e^{i\theta R}$; self-consistency preserves
them into the fixed point rather than generating them from nothing.
\item The Riemannian coupling $\gamma = -1$ produces real-valued fixed
points; $\rho_R$ has only real entries. Classical correlations persist
(the state is entangled, with purity $< 1$) but no quantum coherences
appear.
\item The quantum/classical boundary coincides exactly with the
Lorentzian/Riemannian boundary, confirming the mechanism of
Section~\ref{sec:phase_structure}.
\end{itemize}
This demonstrates that Lorentzian self-consistency can produce quantum
coherences in subsystem density matrices and that Riemannian
self-consistency cannot. Larger models confirm the full algebraic
structure:
\begin{itemize}
\item \textbf{$N = 8$ ($2 \times 2 \times 2$):} All six pair-to-single
subsystem reductions satisfy $\|\mathrm{Tr}_{R_2}\rho_R -
\rho_{R_1}\|_F < 10^{-10}$. The Born rule
$p_k = (\hat{U}\rho_A\hat{U}^\dagger)_{kk}$ holds to $< 10^{-12}$
in computational, Hadamard, and random unitary bases for every
subsystem. The Lorentzian subsystem density matrices have nonzero
imaginary coherences ($|\mathrm{Im}\,\rho_{01}| > 10^{-6}$), while the
Riemannian density matrices are exactly real ($|\mathrm{Im}\,\rho_{01}|
< 10^{-10}$).
\item \textbf{$N = 16$ ($2 \times 2 \times 2 \times 2$):} A
three-level composition chain
$\rho_{ABC} \to \rho_{AB} \to \rho_A = \rho_A^{(\mathrm{direct})}$
passes for all tested paths. Born rule and imaginary coherences
confirmed across five random seeds for the external field~$h$.
\item \textbf{Scaling ($N = 4$--$256$):} The imaginary coherence of
the reduced density matrix $|\mathrm{Im}\,\rho_{01}|$ decays as $\sim
N^{-0.6}$ in the original $\mathrm{dims} = (2,2,\ldots,2)$ setup, and
the reduced-matrix $\mathrm{Im/Re}$ ratio decays as $\sim
d_{\mathrm{env}}^{-0.6}$ when the environment dimension grows. Both are
partial-trace averaging artifacts. The definitive test examines the full
fixed-point wavefunction $\psi^*$ directly: after global phase
alignment, approximately $25\%$ of $\|\psi^*\|$ lies in the imaginary
part for rank-2 bipartitions, rising to $\sim 34\%$ for higher-rank
subsystems, \textbf{stable from $d_{\mathrm{env}} = 2$ through $16$}
with contracting variance ($\pm 0.003$ at $N = 256$). The Riemannian
$\psi^*$ is exactly real. Since the curvature map depends on $\psi$
only through $|\psi_\sigma|^2$, both fixed points satisfy the same
magnitude equation: stripping phases from the Lorentzian $\psi^*$
recovers the Riemannian fixed point by construction. The $68\%$
entropy excess $S_{\mathrm{Lor}} / S_{\mathrm{Riem}} \approx 1.68$
is therefore entirely a phase effect
(Lemma~\ref{lem:entropy_excess}). This ratio is stable from
$N = 4$ to $256$ across all scaling regimes tested.
\end{itemize}
\textbf{Caveat:} the Born rule result is kinematic. Once the
self-consistency map produces $L^2$-normalized complex amplitudes
$\psi^*$, the Born rule $p_k = \mathrm{Tr}(\rho\, P_k)$ follows
automatically from the partial trace construction. What the model shows
is that self-consistency generates the correct \emph{input}---complex
amplitudes with the right normalization and phase structure---not that it
explains why measurement outcomes follow the Born rule. The measurement
problem remains open.

The toy model confirms the algebraic structure of quantum subsystems
(composition, Born rule, imaginary coherences) and provides numerical evidence
that the mechanism survives scaling. The Lorentzian $\psi^*$ retains
its full complex character at every system size tested; only the
reduced-density-matrix representation of this structure fades with
environment dimension, while the entanglement entropy retains its full
$68\%$ excess. The co-emergence chain as a whole remains Conjecture.

%----------------------------------------------------------------------
\subsection{What emergence means here}
\label{sec:emergence_meaning}
%----------------------------------------------------------------------

If $\mu^*$ must be complex for quantum statistics to emerge, then the
claim of the framework must be stated precisely. The topology provides
the configuration space $\mathrm{Sm}(\mathcal{M})$. The self-consistency
condition provides the dynamics (the map $\mathcal{F}$). But the complex
phase comes from the Lorentzian action evaluated on each smooth
structure. The framework derives local Hilbert space structure from
global phase structure via marginalization---the Born rule, the density
matrix, the subsystem-dependence of quantum mechanics are all emergent.
What is not derived from topology alone is the phase itself. This is a
weaker claim than ``quantum mechanics emerges from topology'' but
stronger than ``quantum mechanics is assumed'': the local quantum
structure is genuinely emergent even though the global phase structure is
not.

The critical question is therefore two-part: (1) does the
self-consistency condition $\mathcal{F}(\mu^*) = \mu^*$ naturally
produce complex-valued fixed points even when the input space is purely
topological, or must the phases be imposed by hand? And (2) if $\mu^*$
is complex, does the marginal over a bipartition have the algebraic
structure of a density matrix? Both questions now have affirmative
answers in finite models (Section~\ref{sec:toy_model}): the Lorentzian
self-consistency map produces complex fixed points with $\sim 25$--$34\%$
imaginary content (depending on subsystem rank), stable through $N = 256$; the marginals satisfy
composition and the Born rule, with imaginary coherences in subsystem
density matrices that are absent in the Riemannian case. The phases
restructure subsystem eigenvalue spectra in a way that persists as
entropy excess even when the reduced-matrix off-diagonal phases are
averaged away by the partial trace
(Section~\ref{sec:phase_structure}). A negative answer to~(1) at any $N$
would mean the framework reorganizes quantum mechanics but does not
derive it, and this paper should be read accordingly.

%======================================================================
\section{The co-emergence thesis}
\label{sec:coemergence}
%======================================================================

The results of Section~\ref{sec:phase} demonstrate that Lorentzian
self-consistency produces quantum structure absent in the Riemannian
case. We now propose that this mechanism extends to a co-emergence of
mass, signature, time, and Hilbert space.

Four concepts that are usually treated as independent---mass, Lorentzian
signature, local time, and local Hilbert space---form a closed logical
chain. Each link has a stated rigor level and citation. The chain reads
as a causal sequence; it is not. The four concepts are simultaneous
aspects of the self-consistent block (see Remark~\ref{rem:temporal}
below).

%----------------------------------------------------------------------
\subsection{Mass requires Lorentzian signature}
\label{sec:mass_signature}
%----------------------------------------------------------------------

\emph{(Rigorous.)}

\medskip

\noindent The distinction between physical mass and the parameter $m^2$
is categorical. On a Lorentzian manifold, the dispersion relation
$E^2 = \mathbf{p}^2 + m^2$ defines a mass shell $p_\mu p^\mu = -m^2$---a
hyperboloid in momentum space. Massive particles are irreducible
representations of the Poincar\'e group $\mathrm{ISO}(3,1)$, classified
by mass and spin (Wigner~\cite{wigner}). The mass shell, the dispersion
relation, and the particle classification all require the $(-,+,+,+)$
split.

On a Riemannian manifold, the parameter $m^2$ appears in the field
equation $(\Delta + m^2)\phi = 0$ but does not define a mass shell, a
dispersion relation, or propagating degrees of freedom. The wave operator
$\Delta + m^2$ is elliptic: its solutions decay exponentially rather than
propagate. The Euclidean group $\mathrm{ISO}(4)$ has no analogous
classification of particles by mass and spin. The parameter $m^2$ exists;
the physics of mass does not.

%----------------------------------------------------------------------
\subsection{Lorentzian signature provides local time}
\label{sec:signature_time}
%----------------------------------------------------------------------

\emph{(Rigorous; Geroch~\cite{geroch}, Hawking--Ellis~\cite{hawking_ellis}.)}

\medskip

\noindent Lorentzian signature distinguishes timelike, spacelike, and null
directions. This trichotomy is the physical content of the $(-,+,+,+)$
sign pattern: it is what makes the metric a causal structure and not merely
a distance function. Massive particles follow timelike geodesics, along
which proper time
\begin{equation}
d\tau^2 = -g_{\mu\nu}\,dx^\mu dx^\nu > 0
\label{eq:propertime}
\end{equation}
is well-defined and positive. Massless particles follow null geodesics,
along which $d\tau = 0$.

General relativity prohibits universal time: clocks at different locations
in a gravitational field tick at different rates. But mass provides
\emph{local} time: worldline-dependent clocks carried by massive objects.
Without hyperbolicity, there is no causal structure, no
timelike/spacelike distinction, and no proper time for any object.

%----------------------------------------------------------------------
\subsection{Local time enables local Hilbert space}
\label{sec:time_hilbert}
%----------------------------------------------------------------------

\emph{(Sketch.)}

\medskip

\noindent The Schr\"odinger equation $i\hbar\,\partial|\psi\rangle/\partial t
= H|\psi\rangle$ requires a time parameter $t$. The inner product
$\langle\psi|\phi\rangle$ is defined on a spatial slice at fixed $t$.
Unitarity requires that the inner product be preserved under evolution in
$t$. If there is no universal time (Section~\ref{sec:signature_time}),
there is no universal Hilbert space---the content of
Axiom~\ref{ax:nohilbert}.

But if mass provides local time, a local Hilbert space exists for massive
subsystems that carry their own clocks. The Page--Wootters
mechanism~\cite{page_wootters} makes this precise: apparent time evolution
is recovered from a globally static state by conditioning on a clock
subsystem. The clock must be massive---it must experience proper time along
a timelike worldline. Without mass, the conditioning has no temporal
reference and the mechanism fails.

This is how Axiom~\ref{ax:nohilbert} is realized: Hilbert space is not
fundamental but is derived via the Page--Wootters mechanism from mass and
Lorentzian signature.

\begin{remark}[Acknowledged gap]
\label{rem:pw_gap}
The Page--Wootters mechanism as standardly formulated presupposes a total
Hilbert space $\mathcal{H} = \mathcal{H}_{\mathrm{clock}} \otimes
\mathcal{H}_{\mathrm{system}}$ and a global Hamiltonian constraint
$\hat{H}|\Psi\rangle = 0$. This conflicts with Axiom~\ref{ax:nohilbert}:
the total Hilbert space is assumed, not derived. A fully consistent
treatment requires reformulating the Page--Wootters mechanism without a
fundamental Hilbert space---using instead the marginal of the
smooth-structure measure $\mu^*$ (Section~\ref{sec:level3}). This
reformulation is an open problem (Section~\ref{sec:open}).
\end{remark}

%----------------------------------------------------------------------
\subsection{Local Hilbert space closes the loop}
\label{sec:hilbert_mass}
%----------------------------------------------------------------------

\emph{(Sketch.)}

\medskip

\noindent The Hilbert space for massive quantum fields produces the
quantum stress-energy $\langle\hat{T}_{\mu\nu}\rangle$ that sources
gravity via the semiclassical Einstein equation
\begin{equation}
G_{\mu\nu}[g] = 8\pi G\,\langle\hat{T}_{\mu\nu}\rangle_{g,\mathrm{ren}}.
\label{eq:sce}
\end{equation}
The companion paper~\cite{fixed_point} establishes that eq.~\eqref{eq:sce}
has a self-consistent solution (Rigorous for massive fields with
$m \ll M_P$). Without mass, the only stress-energy source is the conformal
trace anomaly, which drives the Starobinsky de Sitter
solution~\cite{starobinsky}. This solution is self-consistent at the
metric level (Section~\ref{sec:level2}), but it supports no massive
particles, no local clocks, no local Hilbert space, and no effective
quantum mechanics for subsystems. It is a geometry in equilibrium with
vacuum fluctuations---consistent but physically uninhabited.

Mass makes the self-consistent geometry physically inhabited: it provides
the matter content that sources gravity, the timelike worldlines that
define local time, and the clock subsystems that enable local quantum
mechanics. The loop closes:
\begin{equation}
\text{mass} \;\longleftrightarrow\; \text{Lorentzian signature}
\;\longleftrightarrow\; \text{local time}
\;\longleftrightarrow\; \text{local Hilbert space}
\;\longleftrightarrow\; \langle\hat{T}_{\mu\nu}\rangle
\;\longleftrightarrow\; \text{gravity}
\;\longleftrightarrow\; \text{mass}
\label{eq:co_emergence}
\end{equation}
In the framework of Axiom~\ref{ax:timeless}, this is not a circularity
but a self-consistency constraint: the block must satisfy all these
conditions simultaneously.

%----------------------------------------------------------------------
\subsection{The formal statement}
\label{sec:formal_coemergence}
%----------------------------------------------------------------------

\begin{conjecture}[Co-emergence]
\label{conj:coemergence}
Cross-level self-consistency of the four-level hierarchy
(Definition~\ref{def:hierarchy}) has nontrivial solutions---solutions
supporting effective quantum mechanics for subsystems---only for
Lorentzian metric signature with massive field content. Riemannian
and/or massless solutions exist at the metric level but are physically
degenerate: they do not support the Level~3 structure.
\end{conjecture}

The evidence for Conjecture~\ref{conj:coemergence} is summarized in
Table~\ref{tab:signature}. The semiclassical Einstein equation (Level~2)
is completely signature-blind: self-consistent fixed points exist for
all four combinations of signature and field content. The selection
operates at Level~3, where only the Lorentzian-massive combination
supports local time, local Hilbert space, and effective quantum
mechanics.

\begin{table}[h]
\centering
\begin{tabular}{lcccc}
\toprule
& \multicolumn{2}{c}{Massive fields} & \multicolumn{2}{c}{Massless conformal} \\
\cmidrule(lr){2-3} \cmidrule(lr){4-5}
& Level~2 & Level~3 & Level~2 & Level~3 \\
\midrule
Lorentzian & $\checkmark$ (Rigorous) & $\checkmark$ (Sketch) & $\checkmark$ (Rigorous) & $\times$ \\
Riemannian & $\checkmark$ (Sketch) & $\times$ & $\checkmark$ (Rigorous) & $\times$ \\
\bottomrule
\end{tabular}
\caption{Signature selection operates at the cross-level constraint. Level~2
(the SCE) is signature-blind for both massive and massless fields. Only
Lorentzian signature with massive field content supports Level~3.
The Lorentzian massive Level~2 entry is Rigorous (Banach contraction,
$\kappa \sim (m/M_P)^2$, companion paper~\cite{fixed_point}); the
Riemannian massive entry is Sketch (same contraction argument applies to
elliptic operators); the massless entries use the Starobinsky de Sitter /
round $S^4$ solutions~\cite{starobinsky,bunch_davies}.}
\label{tab:signature}
\end{table}

\begin{remark}[The chain is not causal]
\label{rem:temporal}
The co-emergence chain~\eqref{eq:co_emergence} reads as a causal
sequence: mass ``produces'' signature, signature ``provides'' time, time
``enables'' Hilbert space. This temporal reading must be resisted. The
block does not first acquire mass, then a signature, then time, then
quantum mechanics. It satisfies all constraints at all levels
simultaneously. The chain is our explanation of the constraint structure,
expressed in the only order available to linear text, not a description
of how the block came to be.
\end{remark}

%----------------------------------------------------------------------
\subsection{Weakening the mass assumption}
\label{sec:mass_generation}
%----------------------------------------------------------------------

Conjecture~\ref{conj:coemergence} conditions on the existence of massive
field content. This condition can be weakened. The Starobinsky de Sitter
fixed point~\cite{starobinsky} has nonzero scalar curvature
$R = 12H_0^2$, driven by the conformal trace anomaly even for massless
fields. On this background, a scalar field with non-conformal curvature
coupling $\xi \neq 1/6$ acquires an effective mass
\begin{equation}
m_{\mathrm{eff}}^2 = \left(\xi - \tfrac{1}{6}\right) R.
\label{eq:meff}
\end{equation}
For $\xi > 1/6$, this is a genuine positive mass-squared; the field has
a mass shell $p_\mu p^\mu = -m_{\mathrm{eff}}^2$ in the subhorizon WKB
regime, timelike worldlines with proper time, and the capacity to serve
as a Page--Wootters clock. For the conformally coupled case $\xi = 1/6$,
the effective mass vanishes identically---the field remains exactly
massless. This is consistent with the photon, which is conformally
coupled in four dimensions.

The effective mass~\eqref{eq:meff} is the same number on both Lorentzian
de Sitter and Riemannian $S^4$. But on the Riemannian side it is a
correlation length parameter, not physical mass: there is no mass shell,
no dispersion relation, no propagation, and no local clocks. The
signature selectivity of the co-emergence thesis survives with
$m_{\mathrm{eff}}$ in place of bare mass.

\begin{conjecture}[Self-consistent mass generation]
\label{conj:mass_generation}
\emph{(Sketch.)}
Let $(\mathcal{M}, g)$ be a Lorentzian 4-manifold carrying quantum
fields with at least one species having non-conformal curvature coupling
($\xi \neq 1/6$). Then the self-consistent solution of the semiclassical
Einstein equation has $R \neq 0$, and each non-conformal field acquires
effective mass $m_{\mathrm{eff}}^2 = (\xi - 1/6)R$. This effective mass
is sufficient to activate the co-emergence chain: it defines a mass
shell, provides local proper time along timelike worldlines, enables
local Hilbert space via the Page--Wootters mechanism, and sources the
stress-energy that maintains the self-consistent curvature. Conformal
fields ($\xi = 1/6$) remain exactly massless.
\end{conjecture}

Conjecture~\ref{conj:mass_generation} weakens the mass assumption in the
co-emergence thesis from ``massive fields exist'' to ``non-conformally
coupled fields exist.'' The effective mass is generated by the
self-consistent curvature, not assumed as input. The Rigorous components
are the Starobinsky fixed point and eq.~\eqref{eq:meff} (standard
curved-spacetime QFT). The Sketch component is the closure of the
co-emergence chain with $m_{\mathrm{eff}}$ in place of bare mass.

What the conjecture does not achieve: it does not derive the field
content, the curvature coupling constants, the mass spectrum, or fermion
masses (which require additional mechanisms such as Yukawa couplings).
The mass gap origin is partially answered---effective mass is generated,
not assumed---but the residual assumption ($\xi \neq 1/6$ for at least
one species) deserves scrutiny: how strong is it?

\begin{remark}[The residual assumption is weak]
\label{rem:residual_weak}
Conjecture~\ref{conj:mass_generation} does not require the Higgs
field specifically or any particular fundamental scalar. It requires
that at least one bosonic degree of freedom---fundamental or
composite---has $\xi \neq 1/6$. Every scalar field on a curved
background has a curvature coupling $\xi$; the conformal value
$\xi = 1/6$ is a single point in the parameter space of each species.
For the co-emergence chain to fail at this step, \emph{every} scalar
degree of freedom in the effective field theory would need to sit at
exactly the conformal value, with no known symmetry enforcing this
condition.

The argument extends beyond fundamental fields. Effective scalar
degrees of freedom---condensates, order parameters, composite
operators---also couple to curvature, and their effective $\xi$ values
are not protected by conformal invariance. The co-emergence chain
activates as soon as any one of these has $\xi \neq 1/6$. Universal
conformal coupling across all bosonic degrees of freedom would itself
be an extraordinary and unexplained fine-tuning.
\end{remark}

\begin{remark}[Why $\xi = 1/6$ is non-generic]
\label{rem:xi_running}
The remaining assumption can be further motivated by
renormalization group considerations. In the Standard Model at one loop,
the $\beta$-function for the non-minimal coupling
is~\cite{markkanen_sm_curved}
\begin{equation}
16\pi^2 \beta_\xi = \left(\xi - \tfrac{1}{6}\right)
\left[12\lambda + 2Y_2 - \tfrac{3}{2}(g')^2 - \tfrac{9}{2}g^2\right],
\label{eq:beta_xi}
\end{equation}
where $\lambda$ is the Higgs self-coupling, $Y_2 = 3(y_t^2 + \cdots)$
is the Yukawa trace, and $g$, $g'$ are the electroweak gauge couplings.
The entire expression is proportional to $(\xi - 1/6)$: conformal
coupling is an exact one-loop fixed point, even with the full Standard
Model field content.

This fixed point is a UV repeller. The bracket in
eq.~\eqref{eq:beta_xi} evaluates to $\approx +4.8$ at the electroweak
scale (dominated by the top Yukawa contribution $2Y_2 \approx 5.3$).
Positive bracket means $\beta_\xi > 0$ for $\xi > 1/6$ and
$\beta_\xi < 0$ for $\xi < 1/6$: the flow pushes $\xi$ toward $1/6$ in
the infrared and away from $1/6$ in the ultraviolet. At the high
curvatures where the trace anomaly operates, any deviation
$|\xi - 1/6| > 0$ is amplified by the RG flow.

Whether $\xi = 1/6$ is destabilized at two loops and beyond---i.e.,
whether the fixed point itself shifts---is an open question. An explicit
two-loop $\beta_\xi$ computation in curved spacetime has not been
published for the Standard Model. If the fixed point shifts at higher
orders, then $\xi = 1/6$ is not merely non-generic but
\emph{inaccessible}, and the mass generation conjecture closes without
any assumption about field content. If it persists to all orders, then
$\xi = 1/6$ remains a measure-zero fixed point that requires
fine-tuning to reach.

In either case, the argument that non-conformal coupling is generic
does not require RG instability. No known symmetry of the Standard Model
mandates $\xi = 1/6$ for the Higgs field. Conformal invariance is
anomalous in four-dimensional curved spacetime---the trace anomaly that
drives the Starobinsky solution is itself a conformal anomaly. The
symmetry that would protect $\xi = 1/6$ is the same symmetry whose
breaking generates the curvature in the first place. This is a
naturality argument, not a derivation: it establishes that
$\xi \neq 1/6$ is the generic expectation, but does not exclude the
possibility of exact conformal coupling.
\end{remark}

%======================================================================
\section{The self-consistency hierarchy}
\label{sec:hierarchy}
%======================================================================

We propose a four-level hierarchy of self-consistency constraints. The
levels are organized by \emph{coarse-graining}: Level~0 is the coarsest
description, Level~3 the finest accessible to a subsystem. The levels
are not temporal stages. They are simultaneous aspects of the same
four-dimensional block, described at different resolutions. A valid block
must satisfy the constraints at every level, and these constraints
interlock.

\begin{definition}[The four-level hierarchy]
\label{def:hierarchy}
\quad
\begin{itemize}
\item[\emph{Level 0.}] \emph{Topological.} The fundamental object is a
piecewise-linear (PL) 4-manifold $\mathcal{M}$, equipped with no
additional structure.

\item[\emph{Level 1.}] \emph{Smooth.} Given the topological type of
$\mathcal{M}$, the space of compatible smooth structures
$\mathrm{Sm}(\mathcal{M})$ is the configuration space of the theory. A
self-consistency measure $\mu^*$ on $\mathrm{Sm}(\mathcal{M})$ replaces
both the wavefunction and the path integral measure of standard
formulations.

\item[\emph{Level 2.}] \emph{Metric.} Each smooth structure constrains
the space of compatible metrics. The semiclassical Einstein
equation~\eqref{eq:sce} is the fixed-point condition that selects the
metric at this level; its existence is established in the companion
paper~\cite{fixed_point}.

\item[\emph{Level 3.}] \emph{Effective quantum mechanics.} Marginalizing
the smooth-structure measure $\mu^*$ over the complement of a subregion
$\mathcal{R} \subset \mathcal{M}$ produces an effective density matrix
$\rho_\mathcal{R}$ on a local Hilbert space
$\mathcal{H}_\mathcal{R}$. This level functions only for Lorentzian
signature with massive field content (Section~\ref{sec:coemergence}).
\end{itemize}
\end{definition}

The remainder of this section develops each level and specifies the
constraints that link them.

%----------------------------------------------------------------------
\subsection{Level 0: The topological constraint}
\label{sec:level0}
%----------------------------------------------------------------------

The manifold $\mathcal{M}$ must satisfy two requirements simultaneously:
it must admit a Lorentzian metric, and it must support nontrivial smooth
structures. These requirements are in tension.

A closed manifold admits a Lorentzian metric if and only if its Euler
characteristic vanishes: $\chi(\mathcal{M}) = 0$~\cite{geroch}. For
closed simply-connected 4-manifolds---the class where Freedman's
classification theorem~\cite{freedman} is strongest---the Euler
characteristic is $\chi = 2 + b_2$, where $b_2 = \mathrm{rank}(Q_\mathcal{M})$
is the second Betti number. Whenever the intersection form $Q_\mathcal{M}$ is
nontrivial ($b_2 \geq 1$), we have $\chi \geq 3$. The manifold class
where the Freedman--Donaldson machinery provides the richest structure
\emph{cannot carry a Lorentzian metric}.

This obstruction forces the Level~0 manifold away from the closed
simply-connected class. Three alternatives admit both $\chi = 0$ and
nontrivial smooth structures:

\begin{enumerate}
\item \textbf{Non-compact manifolds (exotic $\mathbb{R}^4$).} The
standard $\mathbb{R}^4$ admits uncountably many exotic smooth
structures~\cite{donaldson,taubes}. It has $\chi = 0$ and admits
Lorentzian metrics. However, $H_2(\mathbb{R}^4;\mathbb{Z}) = 0$, so
the intersection form is trivial. The rich topological structure is
carried by the smooth structures themselves, not by the homology.

\item \textbf{Manifolds with boundary (cobordisms).} Compact manifolds
with boundary can have $\chi = 0$ and nontrivial smooth structures. The
boundary data may carry topological information that compensates for the
loss of the closed-manifold intersection form.

\item \textbf{Non-simply-connected manifolds.} Manifolds with nontrivial
fundamental group (e.g., $T^4$) can have $\chi = 0$ and admit Lorentzian
metrics. The theory of exotic structures on non-simply-connected
4-manifolds is less developed than for simply-connected ones.
\end{enumerate}

Which class of 4-manifolds serves as the correct Level~0 input is an
open problem (Section~\ref{sec:open}). The co-emergence thesis does not
depend on the resolution: it requires only that $\mathcal{M}$ admits a
Lorentzian metric and supports enough smooth structures for Level~1 to
be nontrivial.

%----------------------------------------------------------------------
\subsection{Level 1: The smooth-structure measure}
\label{sec:level1}
%----------------------------------------------------------------------

The most remarkable fact in four-dimensional topology is the following:

\begin{theorem}[Donaldson, 1983; Taubes, 1987~\cite{donaldson,taubes}]
\emph{(Rigorous.)}
\label{thm:exotic}
The smooth manifold $\mathbb{R}^4$ admits uncountably many mutually
non-diffeomorphic smooth structures. In every other dimension $n \neq 4$,
$\mathbb{R}^n$ admits a unique smooth structure up to diffeomorphism.
\end{theorem}

The space $\mathrm{Sm}(\mathcal{M})$ of smooth structures compatible
with the PL type of a four-dimensional $\mathcal{M}$ is generically an
uncountable set. In any other dimension, it is a finite set (often a
single point).

\begin{definition}[The smooth-structure measure]
\emph{(Conjecture.)}
\label{def:mu}
The self-consistency measure $\mu^*$ on $\mathrm{Sm}(\mathcal{M})$ is
the unique measure satisfying the self-consistency condition
$\mathcal{F}(\mu^*) = \mu^*$, where $\mathcal{F}$ is the map defined
by: given a trial measure $\mu$ on $\mathrm{Sm}(\mathcal{M})$, compute
the induced metric and matter content at each smooth structure (via
Level~2), compute the resulting stress-energy, solve the Einstein
constraint for a new geometry, and update the weights on
$\mathrm{Sm}(\mathcal{M})$ accordingly.
\end{definition}

\begin{remark}[Relation to the path integral]
The smooth-structure measure $\mu^*$ is an analogue of the gravitational
path integral, but with a structural difference. The standard
gravitational path integral integrates over metrics---a continuous
infinite-dimensional space with an undefined measure. The measure $\mu^*$
integrates over smooth structures on a fixed topological type. While
$\mathrm{Sm}(\mathcal{M})$ may be uncountable, its structure as a moduli
space of gauge-equivalence classes is better controlled than the space of
all metrics. Donaldson and Seiberg--Witten theory provide invariants of
smooth structures that can serve as weights in $\mu^*$, even if the full
measure is not yet defined.
\end{remark}

%----------------------------------------------------------------------
\subsection{Level 2: The metric fixed point}
\label{sec:level2}
%----------------------------------------------------------------------

Given a smooth structure $\sigma \in \mathrm{Sm}(\mathcal{M})$, the
self-consistency constraint at the metric level is
eq.~\eqref{eq:sce}: the geometry that hosts the quantum state must be
sourced by that state. The companion paper~\cite{fixed_point} establishes
existence of solutions at three levels of rigor:

\begin{enumerate}

\item \emph{Exactly}, in the cosmological case via the Starobinsky
(1980) trace anomaly solution~\cite{starobinsky}, with
$H_0^2 = 180\pi/G|a_2|$ (Rigorous).

\item \emph{Perturbatively}, near any classical solution, via a Banach
contraction with constant $\kappa \sim (m/M_P)^2 \ll 1$ for massive
fields (Rigorous). The extension to massless fields remains open: the
spectral gap that controls the infrared behavior of the response kernel
vanishes at $m = 0$.

\item \emph{Conditionally non-perturbatively}, via the Schauder
fixed-point theorem on globally hyperbolic spacetimes with compact Cauchy
surfaces, under six stated assumptions of which four are established and
two remain open (Sketch).

\end{enumerate}

The Planck scale emerges as the natural validity boundary: when
$R \sim \ell_P^{-2}$, the contraction constant $\kappa \to 1$, the
perturbative expansion breaks down, and the full smooth-structure measure
$\mu^*$ becomes relevant.

\begin{remark}[Level~2 is signature-blind]
\label{rem:signature_blind}
The contraction constant $\kappa \sim (m/M_P)^2$ is identical for
Lorentzian and Riemannian metrics. The Banach contraction establishes
existence of self-consistent solutions regardless of signature. For
massless conformal matter, the Starobinsky de Sitter solution
(Lorentzian) and the round $S^4$ (Riemannian) are both exact
self-consistent fixed points~\cite{starobinsky,bunch_davies}. Level~2
does not distinguish signatures for any field content. This is not a
deficiency---it is informative: it demonstrates that signature selection
must operate at the cross-level constraint
(Conjecture~\ref{conj:coemergence}), not at any single level.
\end{remark}

%----------------------------------------------------------------------
\subsection{Level 3: Effective quantum mechanics}
\label{sec:level3}
%----------------------------------------------------------------------

\begin{conjecture}[Emergence of effective quantum mechanics]
\emph{(Conjecture.)}
\label{conj:qm}
Let $\mathcal{R} \subset \mathcal{M}$ be a subregion of the block. Let
$\mu^*$ be the self-consistency measure on $\mathrm{Sm}(\mathcal{M})$.
The marginal of $\mu^*$ over the smooth structures of the complement
$\mathcal{M} \setminus \mathcal{R}$, restricted to $\mathcal{R}$, defines
an effective measure $\mu^*_\mathcal{R}$ on the local configuration space
of $\mathcal{R}$.

When $\mathcal{R}$ is much larger than the Planck scale but much smaller
than the cosmological scale, and gravity within $\mathcal{R}$ is weak,
$\mu^*_\mathcal{R}$ is approximated by a density matrix
$\rho_\mathcal{R}$ on a Hilbert space $\mathcal{H}_\mathcal{R} =
L^2(\mathrm{Config}(\mathcal{R}), \mu^*_\mathcal{R})$.
\end{conjecture}

The Hilbert space $\mathcal{H}_\mathcal{R}$ is derived, not
postulated: it is the $L^2$ space of the marginal measure.
It is subsystem-dependent: different choices of $\mathcal{R}$ give
different Hilbert spaces. There is no single Hilbert space of the
universe. It is approximate: the density matrix approximation is valid
only in the sub-Planckian, weak-gravity regime.

This level functions only for Lorentzian signature with massive field
content. On a Riemannian manifold, the wave operator $\Delta + m^2$ is
elliptic: there are no propagating modes, no mass shell, no particle
interpretation, and no timelike geodesics to provide the local proper
time required for the Page--Wootters mechanism~\cite{page_wootters}. The
parameter $m^2$ is present in the field equation but is not
\emph{mass}: it generates no dispersion, no propagation, no local
clocks. For massless conformal matter on a Lorentzian background, proper
time exists along timelike geodesics, but without massive clock
subsystems the Page--Wootters conditioning has no temporal reference.
Level~3 collapses.

\begin{remark}[Why marginalization produces quantum statistics]
The argument for Conjecture~\ref{conj:qm} parallels the derivation of the
canonical ensemble from microcanonical statistics: integrating out the
environment imposes linearity on the subsystem. In the present context,
the consistency constraints on the full block, restricted to
$\mathcal{R}$, become approximately linear in the degrees of freedom of
$\mathcal{R}$ when the complement is integrated over, because the
dominant contribution from the complement is the mean-field background.
The analogy to statistical mechanics is structural---in both cases,
integrating out a large environment produces an approximately linear
effective theory---but the mechanism differs. Whether the marginalization
of $\mu^*$ produces the same linearization is the content of
Conjecture~\ref{conj:qm}, not a consequence of the analogy.
\end{remark}

%----------------------------------------------------------------------
\subsection{The interlocking constraint}
\label{sec:interlock}
%----------------------------------------------------------------------

The four levels are not independent. The constraint is that the
self-consistency chain must close:
\begin{equation}
\underbrace{\mathcal{M}}_{\text{Level 0}}
\xrightarrow{\text{smooth}}
\underbrace{\mathrm{Sm}(\mathcal{M})}_{\text{Level 1}}
\xrightarrow{\mu^* = \mathcal{F}(\mu^*)}
\underbrace{g^*}_{\text{Level 2}}
\xrightarrow{\text{Einstein}}
\underbrace{\langle\hat{T}_{\mu\nu}\rangle}_{\text{Level 2}}
\xrightarrow{\text{marginalize}}
\underbrace{\rho_\mathcal{R}}_{\text{Level 3}}
\xrightarrow{\text{consistency}}
\mathcal{M}
\label{eq:chain}
\end{equation}

The arrows represent logical constraints, not temporal sequences. A
block satisfying the hierarchy satisfies all constraints simultaneously.
The co-emergence thesis (Conjecture~\ref{conj:coemergence}) is the claim
that the joint solution space of this chain is restricted to Lorentzian
signature with massive field content: this is the only combination for
which all four levels are simultaneously nontrivial.

\begin{remark}[The chain is an epistemic decomposition]
The levels in~\eqref{eq:chain} are introduced for human
comprehensibility. They correspond to the precision at which we can
currently do mathematics. The block does not ``proceed'' through these
levels. It satisfies the constraint. The iteration suggested by
Axiom~\ref{ax:timeless} is our procedure for computing the fixed point,
not a description of how the universe came to be.
\end{remark}

%======================================================================
\section{Why four dimensions}
\label{sec:dimensions}
%======================================================================

The hierarchy of Section~\ref{sec:hierarchy} operates on a 4-manifold.
Here we argue that the four-dimensionality is not an input but a
consequence: four is the unique dimension where the constraint structure
of Definition~\ref{def:hierarchy} is nontrivial.

\begin{table}[h]
\centering
\begin{tabular}{lccc}
\toprule
Dimension & Smooth structures on $\mathbb{R}^n$ & $h$-cobordism theorem & Level~1 richness \\
\midrule
$n < 4$ & Unique & Holds (Smale, $n \geq 6$; special cases $n = 5$) & Insufficient \\
$n = 4$ & Uncountably many & \textbf{Fails} (Donaldson 1987) & Rich \\
$n > 4$ & Unique ($n \geq 5$) & Holds & Insufficient \\
\bottomrule
\end{tabular}
\caption{Topology by dimension. The failure of the $h$-cobordism theorem in
dimension four is the root of the exotic structure richness.}
\label{tab:dimensions}
\end{table}

\begin{proposition}[Four-dimensional uniqueness]
\emph{(Sketch.)}
\label{prop:4d}
The self-consistency hierarchy of Definition~\ref{def:hierarchy} is trivial
in every dimension other than four:
\begin{enumerate}
\item For $n \neq 4$: $\mathrm{Sm}(\mathcal{M})$ is a finite set
(generically a single point for $\mathbb{R}^n$). The Level~1 measure
$\mu^*$ is a finite sum, not an integral over a continuous moduli space.
The effective quantum mechanics of Conjecture~\ref{conj:qm} collapses to
classical statistics.

\item For $n = 4$: $\mathrm{Sm}(\mathcal{M})$ is uncountably infinite.
The Level~1 measure is a genuine integral over a continuous moduli space.
\end{enumerate}
\end{proposition}

\noindent\emph{Sketch.} Point~(1) uses Theorem~\ref{thm:exotic}: smooth
structures on $\mathbb{R}^n$ are unique for $n \neq 4$. The same
conclusion holds for compact manifolds in dimensions $n \geq 5$ (Smale's
$h$-cobordism theorem; see~\cite{smale}), and for $n \leq 3$ by
classical surface and 3-manifold theory. For $n = 4$,
Theorem~\ref{thm:exotic} gives the uncountable smooth structures.
The identification of the uncountable smooth structure moduli space as
the source of quantum statistics (Conjecture~\ref{conj:qm}) is a
conjecture, not a theorem. \hfill$\square$

\begin{remark}
Proposition~\ref{prop:4d} does not \emph{prove} that the universe is
four-dimensional. It establishes that the framework proposed here is
nontrivial only in dimension four. If the framework is correct,
four-dimensionality follows. Whether the framework is correct is an open
question.
\end{remark}

%======================================================================
\section{Connection to existing formulations}
\label{sec:connections}
%======================================================================

\subsection{The Wheeler--DeWitt equation as a Level-2 projection}

The Wheeler--DeWitt equation $\hat{H}\Psi[\gamma_{ij}] =
0$~\cite{dewitt} is the canonical quantization of the metric-level
constraint. In the hierarchy of Section~\ref{sec:hierarchy}, it is an
approximation to the Level~2 fixed-point condition, obtained by:

\begin{enumerate}
\item Concentrating the smooth-structure measure $\mu^*$ on a single
smooth structure (ignoring the Level~1 sum over exotic structures).
\item Performing a $3+1$ decomposition of the metric (introducing a
preferred foliation).
\item Replacing the Level~2 classical fixed-point condition with its
quantum analogue via Dirac quantization.
\end{enumerate}

Each step introduces additional structure not present in the hierarchy.
Step~(2) is the primary source of Axiom~\ref{ax:timeless}
violation~\cite{kuchar}. Steps~(1) and~(3) are simplifications that are
valid in the sub-Planckian, single-smooth-structure regime.

The Wheeler--DeWitt equation is thus the projection of the full hierarchy
onto a single smooth structure and a single spatial foliation. It
contains the correct physics at leading order in $\kappa \sim (m/M_P)^2$,
but not the global, foliation-independent constraint structure of the
full hierarchy.

\subsection{The semiclassical Einstein equation as Level-2}

The companion paper~\cite{fixed_point} studies eq.~\eqref{eq:sce} as a
fixed-point condition in its own right. Within the present hierarchy,
this is the Level~2 constraint evaluated on a single smooth structure.
The Banach contraction result ($\kappa \ll 1$ for sub-Planckian
curvatures) is the perturbative stability of the fixed point at this
level. The Schauder existence result (conditional on two open conjectures)
is non-perturbative existence.

The hierarchy extends this program: Level~2 existence (established with
caveats in the companion paper) is a necessary condition for the full
hierarchy to have a fixed point, but not sufficient. Level~1 requires
additionally that the measure over smooth structures has a fixed point.
Level~0 requires that the manifold type is self-consistent with the
matter content.

\subsection{The Page--Wootters mechanism at Level~3}

The Page--Wootters mechanism~\cite{page_wootters}---apparent time
evolution of a subsystem recovered from a globally static state via
conditioning on a ``clock'' subsystem---is a Level~3 phenomenon in the
present hierarchy. The globally static state is the smooth-structure
measure $\mu^*$; the conditioning on a subsystem is the marginalization
of Conjecture~\ref{conj:qm}. The apparent time evolution is a pattern
in the correlations between $\mathcal{R}$ and its complement encoded in
$\mu^*$.

The clock subsystem must be massive: it must experience proper time
along a timelike worldline (Section~\ref{sec:time_hilbert}). This
requirement connects Level~3 directly to the co-emergence thesis:
without mass, no clock; without a clock, no conditioning; without
conditioning, no effective time evolution and no local Hilbert space.

The observer-dependence identified by Kucha\v{r}~\cite{kuchar} and
confirmed by Marletto and Vedral~\cite{marletto_vedral}---different
choices of ``clock'' subsystem give different dynamics---is, in the
present hierarchy, a manifestation of the fact that different choices
of $\mathcal{R}$ give different marginal measures $\mu^*_\mathcal{R}$
and hence different effective Hilbert spaces
$\mathcal{H}_\mathcal{R}$. The observer-dependence is a feature:
it reflects the emergent, subsystem-relative nature of Hilbert space.

\subsection{Osterwalder--Schrader reconstruction as flat-space evidence}

The Osterwalder--Schrader (OS) reconstruction
theorem~\cite{osterwalder_schrader} provides independent evidence for the
mass-signature link: if a Euclidean quantum field theory has a mass gap
and satisfies reflection positivity, the theory reconstructs uniquely to
a Lorentzian QFT. The chain is: mass gap $\to$ reflection positivity
$\to$ Lorentzian theory. This is a rigorous result on flat space.

However, the OS theorem cannot serve as the \emph{mechanism} for
signature selection in the hierarchy, for two reasons. First, the OS
axioms require a flat Euclidean background with $E(4)$ invariance and a
global reflection hyperplane---fixed background structure that violates
Axiom~\ref{ax:nobackground}. Second, the reflection hyperplane defines a
preferred direction that becomes time after Wick rotation, which violates
Axiom~\ref{ax:timeless}.

The OS theorem is evidence, not mechanism: it demonstrates that the
connection between mass and Lorentzian signature is a rigorous
mathematical phenomenon, not verbal hand-waving. In the hierarchy, it
applies as a consistency check in the flat-space (weak-gravity) limit
of Level~2. The effective QFT obtained by marginalizing $\mu^*$ should
satisfy reflection positivity---a nontrivial constraint on the
self-consistency measure.

%======================================================================
\section{Open problems}
\label{sec:open}
%======================================================================

We list four priority open problems, in order of importance.

\begin{enumerate}

\item \textbf{The mass gap origin.}
The co-emergence thesis requires at least one massive field species.
Conjecture~\ref{conj:mass_generation} shows that effective mass is
generated self-consistently for non-conformal fields, reducing the
assumption from ``massive fields exist'' to ``non-conformally coupled
fields exist.'' Two questions remain open:
\begin{enumerate}
\item[(a)] \emph{Mass existence.} Can the existence of non-conformally
coupled fields ($\xi \neq 1/6$) be derived from Levels~0--1, or is it
irreducible input? The Asselmeyer-Maluga
program~\cite{asselmeyer_maluga} has explored the possibility that
exotic smooth structures generate matter fields from pure topology, but
its core claims remain unverified by independent groups, and the
program operates in Euclidean signature.
\item[(b)] \emph{Mass spectrum.} Can the specific mass values of the
Standard Model be derived from the hierarchy, or do they require
additional structure (field content, coupling constants) as external
input? The trace anomaly mechanism of
Conjecture~\ref{conj:mass_generation} generates a single mass scale
$\sim H_0$, not a spectrum.
\end{enumerate}
A more radical possibility is that mass is the spectral gap of the
self-consistency map $\mathcal{F}$ on the full space of measures on
$\mathrm{Sm}(\mathcal{M})$. If correct, this identification would
subsume both questions into the spectral theory of $\mathcal{F}$. Whether
this identification has mathematical content is an open problem
requiring explicit computation in toy models.

\item \textbf{Level~3 formalization: quantum vs.\ classical.}
The co-emergence chain requires that marginalizing $\mu^*$ over the
complement of a subregion produces \emph{quantum} statistics---density
matrices with complex amplitudes, interference, and Born rule
probabilities---not merely classical stochastic correlations. This
question is now partially resolved. The finite toy model
(Section~\ref{sec:toy_model}) confirms at $N = 8$ and $16$ that
Lorentzian self-consistency produces the full algebraic structure:
subsystem composition ($\rho_{R_1} = \mathrm{Tr}_{R_2}\,\rho_R$ to
machine precision), Born rule statistics in arbitrary measurement bases,
and imaginary coherences in subsystem density matrices---all absent in
the Riemannian case.
A scaling study shows that the imaginary content of the
\emph{reduced density matrix} decays with both $N$ and $d_{\mathrm{env}}$
($\sim d_{\mathrm{env}}^{-0.6}$), but the full fixed-point wavefunction
$\psi^*$ retains $\sim 25\%$ imaginary norm for rank-2 bipartitions
(rising to $\sim 34\%$ for higher-rank subsystems) at all sizes tested
through $N = 256$, while the Riemannian $\psi^*$ is exactly real. The
Lorentzian and Riemannian fixed points share identical magnitude
profiles (the curvature map depends only on $|\psi_\sigma|^2$), so the
$68\%$ entropy excess
($S_{\mathrm{Lor}}/S_{\mathrm{Riem}} \approx 1.68$) is entirely a
phase effect. The reduced-density-matrix decay is a partial-trace
averaging artifact; the mechanism retains its full complex character
in the underlying state. However, the Born rule
result is kinematic: it follows
automatically from $L^2$-normalization and the partial trace, so what
emerges is the correct \emph{input structure} for quantum mechanics, not
an explanation of the measurement process. The open question shifts from
whether quantum \emph{algebraic structure} emerges (it does) to whether
\emph{dynamical} aspects---unitary time evolution, the measurement
process---can be derived from self-consistency.

Separately, the Page--Wootters mechanism as standardly formulated
presupposes a total Hilbert space (Remark~\ref{rem:pw_gap}). Reformulating
it as an operation on measures rather than Hilbert space vectors, and
showing that the conditioning on a massive clock subsystem requires the
metric to be Lorentzian, would promote the co-emergence thesis from
Conjecture to Sketch.

\item \textbf{The Level~0 manifold class.}
Which 4-manifolds admit both Lorentzian metrics ($\chi = 0$) and
nontrivial smooth structures? The Euler characteristic obstruction
(Section~\ref{sec:level0}) rules out closed simply-connected manifolds
with nontrivial intersection forms. The candidates---exotic
$\mathbb{R}^4$, manifolds with boundary, non-simply-connected
manifolds---each have different mathematical properties. Identifying the
correct class is prerequisite to making Level~0 precise.

\item \textbf{A natural measure on smooth structures.}
Define a measure $\mu$ on $\mathrm{Sm}(\mathcal{M})$ for a 4-manifold
$\mathcal{M}$ that is:
\begin{enumerate}
\item Defined without reference to a metric (to avoid the circularity of
the standard gravitational measure problem).
\item Compatible with the self-consistency condition
$\mathcal{F}(\mu) = \mu$.
\item Invariant under the orientation-preserving diffeomorphism group.
\end{enumerate}
Donaldson--Witten theory and Seiberg--Witten invariants provide natural
weights on smooth structures; assembling these into a well-defined
measure on the full moduli space $\mathrm{Sm}(\mathcal{M})$ is the open
problem.

\end{enumerate}

\medskip

\noindent\textbf{Empirical content.} At energies accessible to current
or foreseeable experiments, the hierarchy's empirical content is
identical to semiclassical gravity. The framework does not predict new
physics at sub-Planckian energies; it reorganizes the conceptual
foundations---the relationships between mass, signature, time, and
Hilbert space---in a way that may point toward new physics at Planck-scale
energies. Level~1 corrections (the smooth-structure sum beyond the
single-metric approximation) are suppressed by powers of $\ell_P$ and
invisible at any accessible scale. The companion
paper~\cite{fixed_point} identifies specific predictions (gravitational
decoherence timescales) that arise from the Level~2 fixed point; these
are predictions of semiclassical gravity, not of the full hierarchy.

\medskip

\noindent\textbf{What would falsify the framework.} The co-emergence
thesis (Conjecture~\ref{conj:coemergence}) is a mathematical claim, not
a directly experimental one. It would be falsified by any of the
following, each of which is a definite mathematical question:
\begin{enumerate}
\item \emph{Riemannian quantum mechanics.} A construction of a
self-consistent Riemannian solution of the hierarchy that supports
effective quantum mechanics for subsystems---i.e., a Riemannian manifold
where the marginal of $\mu^*$ produces density matrices with complex
off-diagonal elements, interference, and Born rule statistics despite
the absence of timelike geodesics. The finite toy model
(Section~\ref{sec:toy_model}) finds no such structure at $N = 4$ or at
any $N$ tested through $1024$; a construction at any $N$ would kill the
thesis.

\item \emph{Classical Lorentzian fixed points.} A proof that for
\emph{all} non-uniform weight functions on a finite configuration
space, the self-consistency map $\mathcal{F}$ with Lorentzian coupling
$\gamma = -1 + i\theta$ admits only fixed points whose partial trace
is a real-valued density matrix. The toy model finds complex coherences
at every $N$ tested. Although the imaginary content of the
\emph{reduced density matrix} decays with environment dimension
($\sim d_{\mathrm{env}}^{-0.6}$), the full fixed-point wavefunction
$\psi^*$ retains $\sim 25$--$34\%$ imaginary norm (after global phase
alignment, depending on subsystem rank) at all sizes tested through
$N = 256$, while the Riemannian
$\psi^*$ is exactly real. The two fixed points share identical
magnitude profiles (the curvature map depends only on
$|\psi_\sigma|^2$), so the entire $68\%$ entropy excess
($S_{\mathrm{Lor}}/S_{\mathrm{Riem}} \approx 1.68$)
originates in the phase structure. A proof that $\psi^*$ becomes
effectively real as $N \to \infty$ would falsify the mechanism.
If Lorentzian self-consistency produces only classical correlations in
the continuum limit, the mechanism of Section~\ref{sec:phase} fails.

\item \emph{Born rule failure.} A demonstration that the subsystem
density matrices produced by marginalization of $\mu^*$ fail to satisfy
the full algebraic structure of quantum mechanics---specifically, that
$\rho_R$ does not compose correctly under further bipartitions
($\rho_{R_1} \neq \mathrm{Tr}_{R_2}\,\rho_R$ for $R = R_1 \cup R_2$),
or that measurement outcomes computed from $\rho_R$ violate the Born
rule. Composition and the Born rule have been confirmed through $N = 16$
(Section~\ref{sec:toy_model}); a failure at larger $N$ or with
non-symmetric tensor product structures would be significant.

\item \emph{Mass shell failure.} A proof that curvature-generated
effective mass (eq.~\eqref{eq:meff}) does not produce a physical mass
shell---i.e., that the WKB approximation fails for $m_{\mathrm{eff}}$
in a way that prevents the Wigner classification from applying and the
co-emergence chain from closing.

\item \emph{Smooth-structure triviality.} A proof that all 4-manifolds
admitting Lorentzian metrics have trivial or finite
$\mathrm{Sm}(\mathcal{M})$, leaving no room for the smooth-structure
sum that defines Level~1. This would reduce the hierarchy to standard
semiclassical gravity without the measure-theoretic superstructure.
\end{enumerate}
The framework is structured so that its central claims can be killed by
computation, not only by experiment. The falsification criteria above
are ordered roughly by accessibility: items~1--3 can be tested in
finite toy models of increasing size; items~4--5 require analytic work
in differential geometry.

%======================================================================
\section*{Acknowledgments}
%======================================================================

% [placeholder]

\begin{thebibliography}{99}

\bibitem{fixed_point}
W.~Regelmann and Claude (Anthropic),
``Fixed-point existence of self-consistent semiclassical gravity,''
companion paper (2026).

\bibitem{dewitt}
B.~S.~DeWitt,
``Quantum theory of gravity. I. The canonical theory,''
\textit{Phys.\ Rev.} \textbf{160}, 1113 (1967).

\bibitem{wheeler}
J.~A.~Wheeler,
``Superspace and the nature of quantum geometrodynamics,''
in \textit{Battelle Rencontres}, eds.\ C.~M.~DeWitt and J.~A.~Wheeler
(Benjamin, New York, 1968), pp.~242--307.

\bibitem{kuchar}
K.~V.~Kucha\v{r},
``Time and interpretations of quantum gravity,''
in \textit{Proc.\ 4th Canadian Conference on General Relativity and
Relativistic Astrophysics}, eds.\ G.~Kunstatter, D.~Vincent, and
J.~Williams (World Scientific, 1992), pp.~211--314.

\bibitem{page_wootters}
D.~N.~Page and W.~K.~Wootters,
``Evolution without evolution: Dynamics described by stationary observables,''
\textit{Phys.\ Rev.\ D} \textbf{27}, 2885 (1983).

\bibitem{wigner}
E.~P.~Wigner,
``On unitary representations of the inhomogeneous Lorentz group,''
\textit{Ann.\ Math.} \textbf{40}, 149--204 (1939).

\bibitem{geroch}
R.~Geroch,
``Spinor structure of space-times in general relativity.\ I,''
\textit{J.\ Math.\ Phys.} \textbf{9}, 1739--1744 (1968).

\bibitem{starobinsky}
A.~A.~Starobinsky,
``A new type of isotropic cosmological models without singularity,''
\textit{Phys.\ Lett.\ B} \textbf{91}, 99--102 (1980).

\bibitem{bunch_davies}
T.~S.~Bunch and P.~C.~W.~Davies,
``Quantum field theory in de Sitter space: Renormalization by point-splitting,''
\textit{Proc.\ R.\ Soc.\ Lond.\ A} \textbf{360}, 117--134 (1978).

\bibitem{freedman}
M.~H.~Freedman,
``The topology of four-dimensional manifolds,''
\textit{J.\ Diff.\ Geom.} \textbf{17}, 357--453 (1982).

\bibitem{donaldson}
S.~K.~Donaldson,
``An application of gauge theory to four-dimensional topology,''
\textit{J.\ Diff.\ Geom.} \textbf{18}, 279--315 (1983).

\bibitem{smale}
S.~Smale,
``On the structure of manifolds,''
\textit{Amer.\ J.\ Math.} \textbf{84}, 387--399 (1962).

\bibitem{taubes}
C.~H.~Taubes,
``Gauge theory on asymptotically periodic $4$-manifolds,''
\textit{J.\ Diff.\ Geom.} \textbf{25}, 363--430 (1987).

\bibitem{hawking_ellis}
S.~W.~Hawking and G.~F.~R.~Ellis,
\textit{The Large Scale Structure of Space-Time}
(Cambridge University Press, 1973).

\bibitem{halliwell}
J.~J.~Halliwell,
``Derivation of the Wheeler--DeWitt equation from a path integral for
minisuperspace models,''
\textit{Phys.\ Rev.\ D} \textbf{38}, 2468 (1988).

\bibitem{banks}
T.~Banks,
``TCP, quantum gravity, the cosmological constant and all that\ldots,''
\textit{Nucl.\ Phys.\ B} \textbf{249}, 332 (1985).

\bibitem{marletto_vedral}
C.~Marletto and V.~Vedral,
``Evolution without evolution and without ambiguities,''
\textit{Phys.\ Rev.\ D} \textbf{95}, 043510 (2017).

\bibitem{osterwalder_schrader}
K.~Osterwalder and R.~Schrader,
``Axioms for Euclidean Green's functions,''
\textit{Commun.\ Math.\ Phys.} \textbf{31}, 83--112 (1973);
``Axioms for Euclidean Green's functions.\ II,''
\textit{Commun.\ Math.\ Phys.} \textbf{42}, 281--305 (1975).

\bibitem{asselmeyer_maluga}
T.~Asselmeyer-Maluga and C.~H.~Brans,
``How to include fermions into general relativity by exotic smoothness,''
\textit{Gen.\ Rel.\ Grav.} \textbf{47}, 30 (2015).

\bibitem{markkanen_sm_curved}
T.~Markkanen, S.~Nurmi, A.~Rajantie, and S.~Stopyra,
``The 1-loop effective potential for the Standard Model in curved
spacetime,''
\textit{J.\ High Energy Phys.} \textbf{06}, 040 (2018).

\end{thebibliography}

\end{document}
